% !TeX root = ../../main.tex

\section{Variabili Multiple}

Cosa succede se in un singolo esperimento ci interessano più caratteristiche contemporaneamente?

\begin{equation}
X, Y : \omega \in \Omega \longrightarrow (X(\omega), Y(\omega)) \in \mathcal{X} \times \mathcal{Y} \subseteq \mathbb{R}^2
\end{equation}

In pratica l'evento base è lo stesso, il risultato è formato da una coppia ordinata. Raggruppiamo più misurazioni dello stesso evento in un'unica entità (una coppia, una terna, ecc...).

Per esempio ipotizziamo di avere un elenco dei residenti italiani ($\Omega$) corredato di dati quali altezza ($X$), peso ($Y$) e età ($Z$).
\\Estraiamo una persona a caso ($\omega \in \Omega$), i possibili risultati sono:

\begin{align*}
(X(\omega), Y(\omega)) &\in \mathcal{X} \times \mathcal{Y} & (Y(\omega), Z(\omega)) &\in \mathcal{Y} \times \mathcal{Z} \\
(X(\omega), Z(\omega)) &\in \mathcal{X} \times \mathcal{Z} & (X(\omega), Y(\omega), Z(\omega)) &\in \mathcal{X} \times \mathcal{Y} \times \mathcal{Z}
\end{align*}

Consideriamo la coppia $(X, Y)$.
\\Come calcoliamo la probabilità che si verifichi una specifica combinazione?

\begin{itemize}
    \item Peschiamo $n$ persone a caso ($n$ esperimenti).
    \item Teniamo traccia di tutti i singoli risultati $\{X(\omega_i), Y(\omega_i)\}_{i=1}^n$.
    \item Indichiamo con $n_{X=x, Y=y} = n_{x,y}(n)$ il numero di volte in cui $X=x$ e $Y=y$ (il numero di volte in cui la persona scelta ha altezza $x$ e peso $y$).
    \\In pratica stiamo dicendo che l'evento $\{X=x\} \cap \{Y=y\}$ occorre $n_{X=x, Y=y}$ volte su $n$ esperimenti.
    \item La frequenza relativa è $\frac{n_{X=x, Y=y}}{n}$.
    \item Per ottenere la probabilità dell'evento $\{X=x\} \cap \{Y=y\}$, ci basta passare al limite per $n \to \infty$:
\end{itemize}

\begin{equation}
P_{X,Y}(x,y) = P(\{X=x\} \cap \{Y=y\}) = \lim_{n \to \infty} \frac{n_{X=x, Y=y}}{n} \quad \forall (x,y) \in \mathcal{X} \times \mathcal{Y}
\end{equation}

$P_{X,Y}(x,y)$ è una tabella dove le righe sono elementi di $X$ e le colonne sono elementi di $Y$. La cella della tabella contiene la probabilità che la coppia $(x,y)$ accada. La somma di tutti i numeri della tabella è 1. I numeri nella tabella sono $|\mathcal{X}| \cdot |\mathcal{Y}|$.

Come ogni PMF:
\begin{itemize}
    \item $P_{X,Y}(x,y) \ge 0$
    \item $\sum_{x \in \mathcal{X}} \sum_{y \in \mathcal{Y}} P_{X,Y}(x,y) = 1$
\end{itemize}

\paragraph{Dimostrazione}
\[
1 = P(\Omega) = P\left(\bigcup_{x \in \mathcal{X}} \bigcup_{y \in \mathcal{Y}} \{X=x, Y=y\}\right) = \sum_{x \in \mathcal{X}} \sum_{y \in \mathcal{Y}} \underbrace{P(\{X=x, Y=y\})}_{P_{X,Y}(x,y)}
\]
Questo è vero perché $\{X=x, Y=y\} \cap \{X=x', Y=y'\} = \emptyset$.

\subsection{Proprietà di Marginalizzazione}

Immagina di avere la tabella altezza-peso. Ipotizziamo di voler calcolare la probabilità di estrarre una persona alta 180 cm.
\\Con la tabella sommo l'intera riga dei 180 cm ovvero: $P(180\text{cm}, 50\text{kg}) + P(180\text{cm}, 60\text{kg}) \dots$

Ovvero:
\begin{equation}
\sum_{y \in \mathcal{Y}} P_{X,Y}(x,y) = P_X(x) \quad \text{e} \quad \sum_{x \in \mathcal{X}} P_{X,Y}(x,y) = P_Y(y)
\end{equation}

\paragraph{Dimostrazione}
\[
\{X=x\} = \{X=x\} \cap \Omega = \{X=x\} \cap \bigcup_{y \in \mathcal{Y}} \{Y=y\} = \bigcup_{y \in \mathcal{Y}} (\{X=x\} \cap \{Y=y\})
\]
Quindi:
\[
P(\{X=x\}) = P\left(\bigcup_{y \in \mathcal{Y}} (\{X=x\} \cap \{Y=y\})\right) = \sum_{y \in \mathcal{Y}} P(\{X=x\} \cap \{Y=y\}) = \sum_{y \in \mathcal{Y}} P_{X,Y}(x,y)
\]
Da cui:
\[
P_X(x) = \sum_{y \in \mathcal{Y}} P_{X,Y}(x,y)
\]

\subsection{Variabili Aleatorie Indipendenti}

Due variabili aleatorie $X$ e $Y$ si dicono indipendenti se e solo se gli eventi $\{X=x\}$ e $\{Y=y\}$ sono indipendenti.

In questo caso specifico (indipendenti) la PMF congiunta si fattorizza nel prodotto delle marginali:
\begin{equation}
P_{X,Y}(x,y) = P(\{X=x\} \cap \{Y=y\}) = P(\{X=x\}) \cdot P(\{Y=y\}) = P_X(x) P_Y(y)
\end{equation}

\textbf{N.B.} Questo è l'unico caso. Se le variabili sono collegate non vale più $P_{X,Y}(x,y) = P_X(x) P_Y(y)$.

\subsubsection{Generalizziamo}

Fino ad ora abbiamo parlato di coppie (2 variabili).
\\Sia $(X_1, \dots, X_m) \in \mathcal{X}_1 \times \dots \times \mathcal{X}_m \subseteq \mathbb{R}^m$ una $m$-upla di variabili aleatorie:

Se hai $m$ eventi, tutti indipendenti fra loro, allora la probabilità congiunta si calcola come:
\begin{equation}
P_{X_1 \dots X_m}(x_1, \dots, x_m) = P(\{X_1=x_1\}, \dots, \{X_m=x_m\}) = P(\{X_1=x_1\}) \cdot  \dots  \cdot P(\{X_m=x_m\})
\end{equation}

Ovvero, raggruppando come produttoria:
\[
\prod_{i=1}^m P(\{X_i=x_i\}) = \prod_{i=1}^m P_{X_i}(x_i)
\]

Quindi:
\begin{equation}
\boxed{P_{X_1 \dots X_m}(x_1, \dots, x_m) = \prod_{i=1}^m P_{X_i}(x_i)}
\end{equation}

\subsubsection{Definizione e Intuizione}

Due variabili aleatorie $X \in \mathcal{X}$ e $Y \in \mathcal{Y}$ sono \textbf{indipendenti} $\iff$ gli eventi $\{X=x\}$ e $\{Y=y\}$ sono indipendenti.

Ovvero se la PMF congiunta si fattorizza nel prodotto delle marginali:
\begin{equation}
p_{X,Y}(x,y) = p_X(x) \cdot p_Y(y) \quad \forall (x,y) \in \mathcal{X} \times \mathcal{Y}
\end{equation}

In termini semplici:
\begin{quote}
    \textit{``La probabilità che due cose accadano insieme è esattamente uguale alla probabilità della prima per la probabilità della seconda.''}
\end{quote}

\paragraph{Esempio: Moneta e Dado}
Immagina di lanciare una moneta (Variabile $X$) e un dado a 6 facce (Variabile $Y$).
\\Voglio calcolare la probabilità che escano contemporaneamente testa ($T$) \textbf{E} il numero 6.

Sappiamo dalle probabilità marginali che:
\begin{itemize}
    \item $p_X(\text{testa}) = \frac{1}{2}$ (50\%)
    \item $p_Y(6) = \frac{1}{6}$ (16,6\%)
\end{itemize}

L'alfabeto congiunto della coppia $(X, Y)$ è l'insieme di tutte le possibili combinazioni:
\[ \mathcal{X} \times \mathcal{Y} = \{(T,1), (T,2), \dots, (T,6), (C,1), (C,2), \dots, (C,6)\} \]

Il numero totale di risultati possibili è $|\mathcal{X}| \cdot |\mathcal{Y}| = 12$. 
Quindi, la probabilità congiunta che si verifichi lo specifico evento è:
\[ p_{X,Y}(\text{testa}, 6) = \frac{1}{12} \]

Le variabili sono indipendenti? Verifichiamolo calcolando il prodotto delle marginali:
\[ p_X(T) \cdot p_Y(6) = \frac{1}{2} \cdot \frac{1}{6} = \frac{1}{12} \]

Il risultato coincide ($\frac{1}{12} = \frac{1}{12}$), quindi la risposta è \textbf{SI}!

